We will prove the above statement using the well-ordering principle. For n >= 0. P(n) be the statements

4(1^3 + 2^3 + ... + n^3) = n^2(n + 1)^2

Assume on the contrary that P(n) does not hold for all n >= 0. Then there is some n such that P(n) does not hold. Define

S = {n ∈ Z>=0 : P(n) is false} is nonempty.
By the WOP, S has a least element, say m. Since k is the least element of S, P(n) holds for all 0 <= n < m.

Case 1: m = 0
When n = 0, the sum is empty, so its value is 0. Therefore,
4(1^3 + 2^3 + ... + 0^3) = 0
and
0^2(0 + 1)^2
Thus P(0) holds, contradicting the fact that m is a counterexample
Therefore, m > 0

Case 2: m > 0
Since m is the smallest counterexample, m - 1 cannot be a counterexample. Therefore P(m - 1) is true.
Consider P(m):
4(1^3 + 2^3 + ... + m^3) = 4(1^3 + 2^3 + ... + (m - 1)^3) + 4m^3 = (m - 1)^2m^2 + 4m^3

(m - 1)^2m^2 + 4m^3 = m^2((m - 1)^2+ 4m)
 = m^2(m^2 - 2m + 1 + 4m)
 = m^2(m^2 + 2m + 1)
 = m^2(m^2 + 1)^2

 Therefore, 4(1^3 + 2^3 + ... + m^3) = m^2(m + 1)^2
 Thus, P(m) is true, contradicting the fact that m is a counterexample

 This is a contradiciton because we show that there are no counterexample, after assuming counterexamples exist

 Therefore, P(n) holds for all n >= 0.
